\chapter{Medical Imaging Quality Assurance and Quality Control}
\label{app:imaging_qc}

Reliable diagnostic imaging depends on continuous quality assurance (QA) and quality control (QC) across acquisition hardware, display workstations, radiation dosimetry, and PACS/DICOM workflows.

\section{Modality-Specific Quality Control}

\begin{longtable}{@{}p{0.20\textwidth}p{0.34\textwidth}p{0.38\textwidth}@{}}
\caption{Core quality control checks across major imaging modalities.}
\label{tab:app_imaging_qc} \\
\toprule
\textbf{Domain} & \textbf{Routine QC Checks} & \textbf{Purpose} \\
\midrule
\endfirsthead
\toprule
\textbf{Domain} & \textbf{Routine QC Checks} & \textbf{Purpose} \\
\midrule
\endhead
Radiography / Fluoroscopy & Tube kVp accuracy, half-value layer, AEC reproducibility, exposure index tracking, fluoroscopy dose-rate limits & Maintain image contrast while preventing excessive patient or staff dose \\
Computed Tomography & Daily water phantom CT number (0 HU) and noise check; periodic $\text{CTDI}_{\text{vol}}$ and spatial-resolution verification & Ensure accurate Hounsfield units and consistent low-contrast detectability \\
MRI & Daily SNR, center frequency, geometric linearity, and coil checks; Zone III/IV safety and quench-vent inspection & Prevent image distortion, signal dropout, and ferromagnetic/thermal incidents \\
Ultrasound & Transducer element dropout, dead-zone, axial/lateral resolution, and depth-calibration phantom checks & Detect probe damage and preserve Doppler/grayscale accuracy \\
Nuclear Medicine / PET & Daily gamma-camera uniformity, energy peaking, dose-calibrator constancy, and PET SUV cross-calibration & Ensure accurate quantitative uptake (SUV) and artifact-free scintigraphy \\
Diagnostic Displays & Luminance calibration to DICOM Grayscale Standard Display Function (GSDF) and ambient-light control & Preserve subtle contrast differences across reading workstations \\
\bottomrule
\end{longtable}

\section{Accreditation and Dose Registries}
Programs such as ACR accreditation, IAEA radiation-protection audits, and institutional dose-tracking registries benchmark facility protocols against diagnostic reference levels (DRLs) and ensure ongoing medical-physics oversight.
